\scanpage{71}
\uline{Corollaire} (cf plus bas) Pour tte partie $3$-libre $J$ de $C$, l'ens $\unsure{\mathscr L|L|\mathcal L}(J)$ des él. de $C-J$ non liés aux él. de $J$, est un \uline{hexagone}, et on trouve ainsi une corr. biun. entre $B_3(C)$ et l'ens $\mathrm{Hex}(C)$ des hexagones dans $C$; l'application inverse associe à un hexagone $H$ l'ens $\mathscr L(H)$ des els de $C-H$ non liés aux él. de $H$.
\diagram{\begin{tikzcd}[ampersand replacement=\&,row sep=1.1em,column sep=.75em]
\& |[label=above:{$\unsure{\delta'_6|J'_6|\xi'_6}$}]|\bullet \arrow[-]{r} \arrow[-]{dl} \arrow[-,dashed]{dd} \arrow[-,dashed]{drr}
\& |[label=above:{$\delta_4$}]|\bullet \arrow[-]{dr} \arrow[-,dashed]{dd} \arrow[-,dashed]{dll} \& \\
|[label=left:{$\delta_5$}]|\bullet \arrow[-]{dr} \arrow[-,dashed]{drr} \& \& \&
|[label=right:{$\delta'_5$}]|\bullet \arrow[-]{dl} \arrow[-,dashed]{dll}\\
\& |[label=below:{$\delta'_4$}]|\bullet \arrow[-]{r} \& |[label=below:{$\delta_6$}]|\bullet \&
\end{tikzcd}}

\uline{NB} La donnée d'un graphe hexagonal ($=$ hexagone combinatoire) équivaut à celle d'un couple \deleted{d'ens} $(\tau,K')$ d'un ens de card $2$ et d'un autre de cardinal $3$ (l'hexagone étant alors $\tau\times K'$, deux sommets étant liés ssi ils \deleted{ne sont pas} ont des \deleted{\unsure{projections|projecteurs}} \inserted{composantes} distinctes dans $\tau$ et dans $K'$, i.e. ssi \unsure{ils proviennent} d'un elt de $\tau\times K'$ qui ils \deleted{forment} est le graphe d'une \unsure{imm.|invers.} $\tau\hookrightarrow K'$). Ceci dit, si $(C,J)$ est défini par $(\tau,K,K')$, l'hexagone contenu dans dans $C$ est défini par $(\tau,K')$, ou mieux $\tau\times K'\subset B=\deleted{\unsure{\{b,b'|\{b\cup b'}}\tau\times(K\amalg K')$ \ldots.

\uline{Parties $4$ libres} Cas typique $J=\{\delta_1,\delta_2,\delta_3,\delta_4\}$; $\mathscr L(J)=\{\delta_5,\delta_6,\delta_{56}\}$. \margin{$\unsure{\widehat{\mathrm{Ens}}_2|\mathrm{Ens}_2}\times\widehat{\mathrm{Ens}}_4$} On sait par l'étude précédente que $J$ est contenu dans une base \uline{unique} $b$, et $b\cap \mathscr L(J)$ est év. de cardinal $2$ (c'est $\{\delta_5,\delta_6\}$, il reste un pt dans le \unsure{complém.|complém} ($\delta_{56}$). Donc $\mathrm{Aut}(C,J)$ (ou à fortiori $\mathrm{Aut}(C,\delta_1,--,\delta_4)$) n'opère pas transitivement sur $\mathscr L(J)$, mais \deleted{$\mathscr L(J)$} a deux orbites $b\cap \mathscr L(J)$ ($=\{\delta_5,\delta_6\}$) et le pt fixe $\{\delta_{56}\}$.
% slanted

Donc $\deleted{\unsure{B|\mathrm{Pl}}}\mathrm{Pl}'(L_5,C)\to\mathrm{Pl}(L_4,C)$ est de degré $2$, et $\mathrm{Pl}''(L_5,C)\to\mathrm{Pl}(L_4,C)$ de degré $1$. Notons que l'ens $E(J)$ des el de $C$ liés à \unsure{un|nul|tt} él. de $J$ est de cardinal $2$ (c'est $\{\delta'_5,\delta'_6\}=\unsure{b'-|\}b'|b'\cap}\sigma_b(J)$ \deleted{\unsure{$\mathscr L(J)\to$}}), et on a une bijection canonique (induite par $\sigma_b$)
\[
E(J)\simeq (b-\!J)=b\cap \mathscr L(J)\subset \mathscr L(J)
\]

\uline{Parties $5$ libres} $1^{\mathrm o})$ Celles qui sont contenues dans une base, sont contenues dans une base unique\inserted{-- ce sont les $\mathscr L(J)$ de cardinal $1$ ($=\{\delta_6\}$), et aussi $E(J)$ ($=\{\delta'_6\}$).}
\margin{$\widehat{\mathrm{Ens}}_5$\quad La cat. des $(C,J)$, $J$ une $5$-libre, équivaut \deleted{au \unsure{i}} à celle des \deleted{couples $(C,\unsure{J|j})$ d'une partie de card. $2$, et d'une autre} ens $J$ de card $5$.}
% slanted
$2^{\mathrm o})$ Pour les autres, on vérifie que ce sont des $\unsure{\frac12|4_2|2}$ étoiles.

\uline{Prop.} L'application qui, à tte paire $\{x,y\}$ de pts non liés de $C$, associe $E(x)\cap E(y)$, est une bijection de l'ens de ces paires avec l'ens des $\frac12$ étoiles (parties $5$-libres ne s'immergeant pas dans une base). \inserted{La bijection inverse est $J\longmapsto E(J)$} Si on tient compte de $\{x,y\}\longleftrightarrow(\text{bibase }B=\{b,b'\},\text{ avec }z\in B/\sigma_B)$,
% archive sheet number: 70
